\section{Introduction}
\label{sec:intro}
Changing four mounting holes to six is a shape-editing problem: the hole pattern must change while the flange body, hole diameter and bolt-circle radius remain consistent. Similar requests remove a bracket flange, add a blind shaft slot or enlarge a bore. These operations require localized changes to three-dimensional geometry, and their effects should remain coherent across views. \Cref{fig:benchcad-reference} illustrates two complex released source/target pairs.

Executable CAD provides a structured representation for this task. A model can change a construction history rather than separately generate each rendered view. Yet executable code is only one level of success. A valid edited solid may have the wrong feature locations; a high global overlap score may hide those errors because the affected region is small. For visual-geometry research, the central question is how to evaluate the requested shape change as well as overall shape agreement.

We implement \ours, an instruction-guided program-editing pipeline with three complementary observations: patch and execution outcomes, agreement with reference solids, and source-relative changes in common-camera projections. A supervised language model reads native CAD and an instruction, emits an original-coordinate patch, and executes through a representation-specific backend. Four-view drawings are derived from the same solid. Visual evaluation therefore concerns the executed geometry; it is not an image-conditioning or drawing-perception experiment.

The primary experiment uses a completed Colab T4 run with a fresh Qwen2.5-Coder-3B model, trained on public paired edits from BenchCAD and CAD-Editor~\cite{benchcad2026,cadeditor2025}. The resulting 39 strict matches among 252 held-out tasks improve on the fresh base's one match, but only 39 of 139 executed outputs meet the strict threshold. We retain this failure gap as a substantive finding. A subsequent visual audit illustrates why global agreement cannot substitute for feature-sensitive evaluation.

Mechanical response provides a further, conditional observation. We apply finite-element checks to illustrative edits only after specifying units, material, fixtures and loads. These post-edit checks measure numerical and target-response consistency under synthetic test scenarios. They neither provide training feedback nor establish manufacturing safety.

\paragraph{Contributions.}
\begin{itemize}
\item A documented two-source executable CAD editing experiment, separating patch applicability, solid execution and reference agreement on every retained held-out task.
\item A source-relative multi-view diagnostic over saved edited solids, with a realistic mounting-hole example that exposes the difference between whole-shape similarity and localized edit agreement.
\item An auditable link between edited geometry, derived vector views and scenario-defined FEM fields, with explicit unavailable outcomes and analysis assumptions.
\end{itemize}
Our contribution is empirical evaluation and reproducibility. Standard supervised adaptation, program patches, projections and FEM are not claimed as individually new. Earlier hardware runs, localized transport pilots and deterministic structural controls are documented separately in the supplement.
